\documentclass[11pt]{article}
\usepackage[margin=1in]{geometry}
\usepackage{amsmath,amssymb}
\usepackage{hyperref}

\title{Codex Notes: Neural Wavefunction Plan for Reproducing arXiv:2511.08560}
\author{codex}
\date{2026-06-12}

\begin{document}
\maketitle

\section{Scope and status}

These notes record an early-stage plan for using neural variational wavefunctions
to reproduce observables studied in Cho, Gabai, Lin, Yeh, and Zheng,
\emph{Bootstrapping Euclidean Two-point Correlators},
\href{https://arxiv.org/abs/2511.08560}{arXiv:2511.08560}.
The paper itself does not use neural wavefunctions.  It derives rigorous
bootstrap bounds on Euclidean two-point correlators and then extracts low-lying
adjoint-sector gaps and matrix elements from those bounds.

Therefore the proposed neural calculation is not a literal reproduction of the
paper's method.  It is an independent variational calculation of the same
spectral data and correlators.  Any result from this route should be described
as approximate unless accompanied by controlled convergence, exact benchmarks,
and independent checks.

\section{Target model}

The target system is the ungauged one-matrix quantum mechanics with Hamiltonian
\begin{equation}
  H = \frac{N}{2}\,{\rm tr}\left(\frac{1}{2}P^2 + V(X)\right),
  \qquad
  V(x) = \frac{1}{2}x^2 + g x^4 ,
\end{equation}
using the paper's convention
\begin{equation}
  {\rm tr}\,O = \frac{1}{N}{\rm Tr}\,O,\qquad
  [X_{ab},P_{cd}] = \frac{i}{N}\delta_{ad}\delta_{bc}.
\end{equation}
The Hilbert space is ungauged, so non-singlet representations of \(U(N)\) are
included.  The ground state is expected to be a singlet, while the correlator
\(\langle \Omega|{\rm tr}\,X(\tau)X(0)|\Omega\rangle\) probes adjoint states.
At large \(N\), the adjoint spectrum is benchmarked by the
Marchesini--Onofri singular integral equation.

\section{Spectral reconstruction of correlators}

The neural goal should be to approximate the ingredients in the spectral
decomposition
\begin{equation}
  G_O(\tau)
  =
  \langle \Omega|O(\tau)O(0)|\Omega\rangle
  =
  \sum_n |\langle \Omega|O|n\rangle|^2
  e^{-(E_n-E_0)\tau},
\end{equation}
where \(|n\rangle\) lies in the representation selected by \(O\).  For
\(O={\rm tr}\,X\) or \(X_{ab}\), this is the adjoint sector.  The deliverables
are therefore:
\begin{enumerate}
  \item \(E_0\) and ground-state diagnostics;
  \item adjoint gaps \(\Delta_n=E_n-E_0\);
  \item transition matrix elements \(|\langle \Omega|O|n\rangle|^2\);
  \item reconstructed correlators with error estimates and comparison to
        bootstrap/Marchesini--Onofri data.
\end{enumerate}

\section{Recommended architecture}

The first serious implementation should not use an unconstrained dense neural
network on the \(N^2\) matrix entries.  It should build in the \(U(N)\) symmetry
from the start.

\subsection{Ground-state singlet ansatz}

For the singlet ground state, work first in eigenvalue variables.  After
absorbing the Vandermonde factor, the singlet sector maps to an antisymmetric
\(N\)-fermion problem.  A suitable ansatz is a Slater--Jastrow/backflow
wavefunction
\begin{equation}
  \Phi_\theta(\lambda_1,\ldots,\lambda_N)
  =
  \det[\varphi_i^\theta(y_j)]\,
  \exp\!\left[
    \sum_i u_\theta(\lambda_i)
    + \sum_{i<j} v_\theta(\lambda_i,\lambda_j)
  \right],
\end{equation}
with backflow coordinates
\begin{equation}
  y_i = \lambda_i + b_\theta(\lambda_i,\{\lambda_j\}_{j\ne i}).
\end{equation}
The one-body orbitals should be initialized to harmonic-oscillator orbitals for
the \(g=0\) benchmark.  This gives the architecture the exact solvable limit as
a natural starting point.

\subsection{Adjoint excited-sector ansatz}

For finite \(N\), an adjoint wavefunction should transform equivariantly,
\begin{equation}
  \Psi_{ab}(UXU^\dagger)=U_{ac}\Psi_{cd}(X)U^\dagger_{db}.
\end{equation}
A conservative equivariant form is
\begin{equation}
  \Psi^{(n)}_{ab}(X)
  =
  \Psi_{0,\theta}(X)\,
  \left[
    \sum_{k=0}^{K} c^{(n)}_k(s_\theta(X))
    \left(X^k-\frac{1}{N}{\rm Tr}(X^k)\mathbf{1}\right)
  \right]_{ab},
\end{equation}
where \(s_\theta(X)\) denotes invariant trace features such as
\({\rm tr}\,X^m\).  This form preserves the adjoint transformation law and
enforces tracelessness.  Parity under \(X\mapsto -X\) should be imposed by
choosing even or odd powers according to the desired \(Z_2\) sector.

For large \(N\), it may be more efficient and more directly comparable to the
paper to solve the Marchesini--Onofri problem with a one-dimensional spectral
or neural basis for \(\psi_n(x)\), using the large-\(N\) eigenvalue density as
input.  That calculation should be treated as a benchmark and diagnostic for
the full neural variational approach.

\section{Training objectives}

Ground-state training should use variational Monte Carlo with the Rayleigh
quotient
\begin{equation}
  E_\theta =
  \frac{\langle \Psi_\theta|H|\Psi_\theta\rangle}
       {\langle \Psi_\theta|\Psi_\theta\rangle},
\end{equation}
estimated with samples from \(|\Psi_\theta|^2\).  Each reported energy must
include sample count, random seed, acceptance rate or effective sample size,
and an uncertainty estimate.

Excited states should be trained by minimizing the Hamiltonian in a symmetry
sector with explicit orthogonality constraints.  In practice this means solving
the generalized eigenvalue problem in a learned basis or adding a rigorously
tracked penalty
\begin{equation}
  {\cal L}_n =
  E_n[\Psi_n]
  +
  \sum_{m<n}\alpha_m
  |\langle \Psi_m|\Psi_n\rangle|^2,
\end{equation}
with overlaps measured from independent Monte Carlo samples.  A learned-basis
generalized eigenvalue solve is preferred once the infrastructure exists,
because it provides clearer orthogonality diagnostics.

\section{Implemented harmonic oscillator benchmark}

As the first executable benchmark, the repository now includes a PyTorch
implementation of the one-dimensional quadratic Hamiltonian
\begin{equation}
  H = -\frac{1}{2}\frac{d^2}{dx^2}
      +\frac{1}{2}\omega^2 x^2 ,
\end{equation}
in units \(\hbar=m=1\).  This example is deliberately simpler than the
finite-\(N\) matrix models.  Its purpose is to verify the wavefunction,
autodifferentiation, quadrature, and Monte Carlo diagnostics against exact
answers before adding matrix degrees of freedom.

\subsection{Ansatz and symmetry}

The implemented ground-state ansatz is a real positive wavefunction written as
\begin{equation}
  \log \psi_\theta(x)
  =
  -\frac{1}{2}\alpha x^2 + f_\theta(x^2),
  \qquad \alpha>0 .
\end{equation}
The Gaussian envelope supplies normalizable tails, while the neural correction
uses the feature \(x^2\).  Therefore
\begin{equation}
  \psi_\theta(x)=\psi_\theta(-x)
\end{equation}
exactly, as required for the ground state of an even potential.  This is an
implementation convenience, not a proof of correctness: it only builds in the
known parity of this benchmark.

\subsection{Energy functional}

Quadrature training minimizes the Rayleigh quotient in gradient form,
\begin{equation}
  E_\theta =
  \frac{
    \int dx\,|\psi_\theta(x)|^2
    \left[
      \frac{1}{2}\left(\partial_x\log\psi_\theta(x)\right)^2
      + V(x)
    \right]
  }{
    \int dx\,|\psi_\theta(x)|^2
  } .
\end{equation}
This is equivalent to the usual kinetic term
\(\frac{1}{2}\int dx\,|\partial_x\psi_\theta|^2\), assuming the boundary term
vanishes.  For the harmonic oscillator and the Gaussian envelope, this
assumption is satisfied on the continuum problem; the finite quadrature grid
must still be checked by increasing the cutoff and grid size.

The local energy is computed diagnostically as
\begin{equation}
  E_L(x)
  =
  -\frac{1}{2}
  \left[
    \partial_x^2\log\psi_\theta(x)
    +
    \left(\partial_x\log\psi_\theta(x)\right)^2
  \right]
  + V(x).
\end{equation}
For an exact eigenstate, \(E_L(x)\) is constant and its variance vanishes.

\subsection{Exact checks}

For the exact ground state
\begin{equation}
  \psi_0(x) \propto \exp\left(-\frac{1}{2}\omega x^2\right),
\end{equation}
the benchmark values are
\begin{equation}
  E_0=\frac{\omega}{2},\qquad
  \langle T\rangle=\langle V\rangle=\frac{\omega}{4},
\end{equation}
\begin{equation}
  \langle x^2\rangle=\frac{1}{2\omega},
  \qquad
  \langle x^4\rangle=\frac{3}{4\omega^2}.
\end{equation}
The exact Euclidean two-point function is
\begin{equation}
  G(\tau)=\langle x(\tau)x(0)\rangle
  =\frac{1}{2\omega}e^{-\omega\tau}.
\end{equation}
The implemented ground-state calculation directly checks \(G(0)=\langle
x^2\rangle\).  Reconstructing nontrivial correlators from learned states
requires excited-state wavefunctions and transition matrix elements, and is
therefore outside this first benchmark.

\subsection{Notebook and script workflow}

The harmonic oscillator notebook is allowed to run training directly because
the problem is one-dimensional, fast, and exactly benchmarked.  This should not
be the default pattern for future matrix-model calculations.  For heavier or
production-quality runs, training should be done by reproducible Python scripts
that print diagnostics and, when appropriate, save structured result files.
Notebooks should then load those saved results, compare them to analytic or
bootstrap benchmarks, and present the interpretation.  This separation keeps
expensive numerical work reproducible and avoids notebook-only artifacts.

\subsection{Krylov time evolution for correlators}

The repository now also includes a Krylov/Lanczos time-evolution utility for
Euclidean correlators.  Given a normalized ground state and an operator \(O\),
define the source state
\begin{equation}
  |\phi\rangle = O|\Omega\rangle .
\end{equation}
The correlator is
\begin{equation}
  G_O(\tau)
  =
  \langle \phi|e^{-\tau(H-E_0)}|\phi\rangle .
\end{equation}
Rather than computing all excited states, the Lanczos algorithm builds the
Krylov space
\begin{equation}
  {\cal K}_m(H,|\phi\rangle)
  =
  {\rm span}\{|\phi\rangle,H|\phi\rangle,\ldots,H^{m-1}|\phi\rangle\}.
\end{equation}
After orthonormalization, the Hamiltonian is represented by a small tridiagonal
matrix \(T_m\), giving the approximation
\begin{equation}
  G_O(\tau)
  \approx
  \|\phi\|^2
  \left[
    e^{-\tau(T_m-E_0)}
  \right]_{00}.
\end{equation}
Diagonalizing \(T_m\) gives Ritz values \(\epsilon_j^{(m)}\) and first-component
weights \(|u_{0j}|^2\), so equivalently
\begin{equation}
  G_O(\tau)
  \approx
  \|\phi\|^2
  \sum_j |u_{0j}|^2
  e^{-\tau(\epsilon_j^{(m)}-E_0)} .
\end{equation}
The values \(\epsilon_j^{(m)}\) are source-sector energy estimates, not a full
diagonalization of the Hamiltonian.  Small or vanishing spectral weights mean
that a Ritz value contributes weakly or not at all to the chosen correlator.
For the harmonic oscillator with \(O=x\), the source \(x|\Omega\rangle\) is
proportional to the first excited state, so the Krylov space closes after one
direction in the continuum theory.  On a finite-difference grid this closure is
approximate and should be checked by grid refinement.

\section{Verification ladder}

\begin{enumerate}
  \item Harmonic oscillator, one degree of freedom: reproduce exact ground and
        excited energies and correlators.
  \item Harmonic one-matrix model at finite \(N\): reproduce the known Slater
        determinant/eigenvalue result and the exact Euclidean correlator.
  \item Quartic stable case \(g>0\): compare against exact diagonalization for
        very small \(N\), where possible.
  \item Large-\(N\) quartic model: compute the eigenvalue density and solve the
        Marchesini--Onofri equation; this is the reference calculation for
        adjoint gaps and matrix elements.
  \item Neural finite-\(N\) runs: extrapolate \(N\to\infty\) and compare gaps,
        matrix elements, and correlators against the bootstrap/MO values.
  \item Near-critical negative coupling: only after the previous checks are
        stable, attempt the paper's near-critical values such as
        \(g=-0.075026\).  Claims here must state the metastability/large-\(N\)
        assumptions clearly.
\end{enumerate}

\section{Repository structure to build later}

A small, auditable codebase should separate:
\begin{itemize}
  \item \texttt{basis/}: eigenvalue, trace, and equivariant adjoint features;
  \item \texttt{ansatz/}: Slater--Jastrow/backflow and adjoint heads;
  \item \texttt{hamiltonian/}: local energy and convention-checked operators;
  \item \texttt{sampler/}: Metropolis or Langevin sampling with diagnostics;
  \item \texttt{estimators/}: energy, overlaps, matrix elements, correlators;
  \item \texttt{benchmarks/}: harmonic oscillator and small-\(N\) checks;
  \item \texttt{notes/}: derivations and convention checks.
\end{itemize}

\section{Open issues}

\begin{itemize}
  \item The phrase ``approximate exactly'' should be avoided in final claims.
        Neural wavefunctions provide systematically improvable approximations,
        not exact states, unless a solvable benchmark is being represented by
        construction.
  \item The finite-\(N\) ungauged adjoint Hilbert space and the large-\(N\)
        Marchesini--Onofri description must not be conflated.
  \item Computing \(G(\tau)\) from a finite number of learned states is reliable
        only over a time window where omitted higher states are controlled.
  \item For negative quartic coupling, finite-\(N\) tunneling makes the model
        unstable.  The calculation must explicitly state when it is targeting
        the intrinsically large-\(N\) metastable/critical problem.
\end{itemize}

\end{document}
